% TOP_PROBLEM: {"id":30007555,"problem_number":"LOCAL-30007555","title":"Identifying choice procedures","category_id":21,"category":{"id":21,"name":"economics","display_name":"Economics","description":"Open economic theory statements, published research agendas, and proposed scoped specifications; question provenance and historical priority are qualified per record.","slug":"economics","order_index":21,"created_at":"2026-10-08T00:00:00Z"},"status":"research_question","external_url":null,"legacy_ids":[],"published":false,"statement":"Characterize when feasible menu and process-data experiments distinguish competing choice procedures, and design cost-constrained experiments with valid confidence sets under model overlap.","economics":{"original_id":44,"field":"Economic theory, equilibrium and social choice","kind":"identification","formalization_basis":"adapted_research_question","source_status":"research_agenda","priority_status":"earliest_found","priority_note":"This is the earliest source in which the relevant agenda was verified during this audit; absolute priority has not been established.","earliest_verified_reference":{"authors":["Geoffroy de Clippel","Kareen Rozen"],"title":"Bounded Rationality in Choice Theory: A Survey","year":2023,"url":"https://geoffroydeclippel.net/Working%20Papers%20PDFs/JEL.pdf","doi":"10.1257/jel.20231592","locator":"§6.4 Searching for common ground; §7 Looking Forward, PDF pp.51–54","evidence_summary":"The revised working paper identifies welfare comparisons, additional process data, and computational or procedural constraints as continuing research directions."},"current_reference":{"authors":["Geoffroy de Clippel","Kareen Rozen"],"title":"Bounded Rationality in Choice Theory: A Survey","year":2023,"url":"https://geoffroydeclippel.net/Working%20Papers%20PDFs/JEL.pdf","doi":"10.1257/jel.20231592","locator":"§6.4 Searching for common ground; §7 Looking Forward, PDF pp.51–54","evidence_summary":"The revised working paper identifies welfare comparisons, additional process data, and computational or procedural constraints as continuing research directions."},"formalization_note":"The classes, experimental design criterion, and separation problem are our scoped adaptation of the survey’s process-data agenda."},"statusQualification":"Published question or research agenda verified at the cited locator. The mathematical specification is adapted for this catalogue and includes additional assumptions; source publication does not establish first-ever priority or certify this exact model remains unsolved. The classes, experimental design criterion, and separation problem are our scoped adaptation of the survey’s process-data agenda.","statusReviewedAt":"2026-10-08"}
% FIRST_POSTED: 2026-10-08
% ENTRY_KIND: statement_only
% !TeX program = lualatex
\documentclass[11pt]{article}
\usepackage{fontspec}
\usepackage[a4paper,margin=25mm]{geometry}
\usepackage{amsmath,amssymb,mathtools}
\usepackage{xurl}
\usepackage[hidelinks]{hyperref}
\newcommand{\sref}[2]{\hyperlink{source:#1:#2}{[S#2]}}
\setlength{\emergencystretch}{3em}
\begin{document}
% problemId:30007555
\section*{Identifying choice procedures}\label{30007555}
\subsection{Definitions and mathematical statement}
Take a finite alternative set \(X\), a collection of menus \(\mathcal B\), and a set \(\mathcal T\) of feasible experimental interventions on presentation, search costs, and decision time. Let \(\mathcal M_1,\ldots,\mathcal M_K\) be specified parametric classes of choice procedures, including limited attention, satisficing, and uncertain preferences. Each pair \((k,\theta)\), where \(\theta\in\Theta_k\), induces a joint law \(P_{k,\theta}^{t,B}\) of chosen alternative and process observations under intervention \(t\) and menu \(B\). Choices alone may leave these classes observationally overlapping. An experiment \(e\) is a probability distribution over \(\mathcal T\times\mathcal B\), and its discrimination criterion is
\[
 J(e)=\inf_{k\ne\ell,\theta\in\Theta_k,\eta\in\Theta_\ell}
 \sum_{t,B}e(t,B)\,D_{\mathrm{KL}}(P_{k,\theta}^{t,B}\Vert P_{\ell,\eta}^{t,B}),
\]
where \(D_{\mathrm{KL}}\) is relative entropy. Identify when feasible interventions yield positive separation, characterize unavoidable equivalence classes, and design experiments maximizing discrimination subject to a stated cost budget. Establish sample requirements for valid model confidence sets under overlap, rather than forcing a single winner. Process measures must have explicitly modeled measurement error, and interventions must not covertly change the economic alternatives. The research target is reliable distinction between economically different procedures and their counterfactual implications, not merely fitting observed choices more closely.

\par\textbf{Formulation and scope.} The classes, experimental design criterion, and separation problem are our scoped adaptation of the survey’s process-data agenda.

\par\textbf{Open-question provenance.} An explicit question or research agenda was verified at the source locator. This does not by itself certify that every later version remains unresolved \sref{30007555}{1}.

\par\textbf{Historical priority.} This is the earliest source in which the relevant agenda was verified during this audit; absolute priority has not been established.

\subsection{Short English statement}
Characterize when feasible menu and process-data experiments distinguish competing choice procedures, and design cost-constrained experiments with valid confidence sets under model overlap.

\subsection{Sources}
\begin{itemize}\raggedright
\item[\textnormal{[S1]}]\hypertarget{source:30007555:1}{} Geoffroy de Clippel, Kareen Rozen. 2023. Bounded Rationality in Choice Theory: A Survey. §6.4 Searching for common ground; §7 Looking Forward, PDF pp.51–54.
\url{https://geoffroydeclippel.net/Working%20Papers%20PDFs/JEL.pdf}.
DOI: 10.1257/jel.20231592.
{\small\textit{Earliest verified question/agenda reference; historical priority qualified below. The revised working paper identifies welfare comparisons, additional process data, and computational or procedural constraints as continuing research directions.}}
\end{itemize}
\end{document}
